% Ampliación aditiva 2026-09-27. Propietarios y verificadores: recibo editorial contiguo.
\section{Composición de acción, área y lectura térmica marcada}
\label{sec:ii-termica-marcada}

La construcción precedente determina un portador graduado, su memoria
energética, las corrientes entre sectores y una referencia espectral marcada.
Esta sección establece cómo actúan esos datos sobre los lectores de energía
e información del artículo. El desarrollo tiene dos niveles demostrativos:
la composición de las secciones de acción, área y horizonte ya construidas,
y una realización termométrica explícita mediante dos funcionales del mismo
registro ampliado. La segunda operación fija un coeficiente térmico dentro
de la regla de lectura que se enuncia y demuestra aquí.

Los antecedentes conservan la secuencia APP--TRIT--TPK: las hojas de suma
y producto retienen residuo y cociente; TRIT conserva régimen y orientación;
TPK transporta ruta, frontera, acarreo y memoria. La conexión nonádica retorna
la fase mientras prolonga el estado enriquecido. Las coordenadas
$\pi_{\HMT},\varphi_{\HMT},e_{\HMT},\alpha_{\HMT}$, la acción
orientada y la incidencia hexada--octada son salidas anteriores de esa misma
estructura discreta conjunta del continuo. En las ecuaciones siguientes,
$\pi,\alpha$ y los demás lectores escalares conservan esa procedencia.

\subsection{Energía de retorno y resolución areal}
\label{subsec:ii-termica-accion-area}

Fijemos una orientación $a$ y la realización radial del artículo. Escribimos
\begin{equation}
 L_\alpha=\frac{5759}{23040}\alpha^{16}L_* ,\qquad
 T_{108}=108t_0=\frac{2\pi L_\alpha}{c},\qquad
 \epsilon_a=\frac{h_a}{T_{108}}=\frac{\hbar_a c}{L_\alpha}.
 \label{eq:ii-ter27-energia}
\end{equation}
El símbolo $L_\alpha$ designa una longitud; el grado energético $L$ del
portador precedente es un operador adimensional distinto. El funcional
orientado y el carácter angular previamente construidos son
\begin{align}
 \gamma&=\frac{\pi AC^*}{180}
 +12\bigl[S_{90}(q_-)-S_{90}(q_+)\bigr]
 -\bigl[S_{120}(q_-)-S_{120}(q_+)\bigr],\nonumber\\
 S_m(q)&=\frac{q^m}{1-q^{3m}},\qquad
 a_0=\frac{1+2\cos(10^\circ)}3\frac{C^*}{6}\frac{\pi}{180}.
 \label{eq:ii-ter27-barbero}
\end{align}
En su sector positivo, el cuanto de área es
$q_A=\gamma a_0L_\alpha^2$. Definamos la energía por unidad de área
$\mathcal T_{I/A}=\epsilon_a/q_A$.

\begin{proposition}[Composición areal, constitutiva y gravitatoria]
\label{prop:ii-ter27-area}
Con el coeficiente radial $G_a\hbar_a=c^3L_\alpha^2$, se cumple
\begin{align}
 \mathcal T_{I/A}&=\frac{c^4}{\gamma a_0G_aL_\alpha},&
 \frac{\mathcal T_{I/A}q_A}{\ln3}
 &=\frac{h_a}{108t_0\ln3},\label{eq:ii-ter27-area-energy}\\
 G_a\,k_B\Theta_{{\rm clk},a}\ln3&=c^4L_\alpha.&
 \label{eq:ii-ter27-grav-clock}
\end{align}
Si $\mathfrak e_a$ denota la carga de la sección constitutiva
$\mathfrak e_a^2=2\alpha h_a/Z_{\rm vac}$, la misma energía por nat es
\begin{equation}
 \frac{\epsilon_a}{\ln3}
 =\frac{Z_{\rm vac}\mathfrak e_a^2}{216\alpha t_0\ln3}.
 \label{eq:ii-ter27-constitutive}
\end{equation}
\end{proposition}
\begin{proof}
Sustituir $q_A$ y \eqref{eq:ii-ter27-energia} da
$\mathcal T_{I/A}=\hbar_ac/(\gamma a_0L_\alpha^3)$.
La identidad radial transforma esta expresión en la primera de
\eqref{eq:ii-ter27-area-energy}; multiplicar por $q_A$ prueba la segunda.
La relación térmica ya establecida es
$k_B\Theta_{{\rm clk},a}\ln3=\epsilon_a$.
Su producto por $G_a$ cancela $\hbar_a$ y prueba
\eqref{eq:ii-ter27-grav-clock}. Finalmente,
$h_a=Z_{\rm vac}\mathfrak e_a^2/(2\alpha)$ en
\eqref{eq:ii-ter27-energia} produce \eqref{eq:ii-ter27-constitutive}.
\end{proof}

Barbero determina la resolución areal que interviene en la tensión. El
producto de ambas lecturas recupera la energía de retorno. Esta cancelación
explica la compatibilidad entre las formulaciones geométrica, constitutiva
y térmica: las tres transportan la misma sección de acción.

\subsection{Información de procedencia y respuesta de horizonte}
\label{subsec:ii-termica-horizonte}

Sobre los treinta y seis factores tríticos de borde conservamos
$N=N_{90}+N_{120}$ y $D=90N_{90}+120N_{120}$. El área es
$\widehat{\mathcal A}=q_AN$, mientras que el estado sectorial y su
Hamiltoniano modular son
\begin{equation}
 \rho_A=Z_A^{-1}e^{-\Delta D},\qquad
 K_A=-\ln\rho_A=(\ln Z_A)I+\Delta D.
 \label{eq:ii-ter27-modular}
\end{equation}
Aquí $\Delta$ es finito y $Z_A$ la normalización conocida. La incidencia
permite recuperar
\[
 N_{90}=4\widehat{\mathcal A}/q_A-D/30,\qquad
 N_{120}=D/30-3\widehat{\mathcal A}/q_A.
\]
Para obtener $D$ desde $K_A$ se utiliza adicionalmente $\Delta\ne0$.
Las configuraciones $(1,0)$ y $(0,1)$ tienen igual área $q_A$ y costes
modulares $90\Delta$ y $120\Delta$. Su diferencia procede de la
incidencia sectorial que conserva el estado, además de su ocupación total.

La purificación de \ref{prop:ii-area-purificacion-producto} conserva el
estado completo y su bipartición. En una realización de horizonte de radio
$R>0$, la escala energética y el operador de respuesta son
\begin{equation}
 \tau_R=\frac{\hbar_ac}{2\pi R},\qquad
 H_R=\tau_R\Delta D+e_RI,
 \label{eq:ii-ter27-response}
\end{equation}
con origen $e_R$ registrado. Si $\delta E_R=\Tr[(\sigma-\rho_A)H_R]$
y $\delta s=s(\sigma)-s(\rho_A)$, el origen cancela en el incremento.

\begin{theorem}[Balance finito y conversión entre secciones térmicas]
\label{thm:ii-ter27-horizon}
Para todo estado $\sigma$ del portador sectorial,
\begin{equation}
 \delta E_R-\tau_R\delta s
 =\tau_R D(\sigma\Vert\rho_A)\ge0.
 \label{eq:ii-ter27-free-area}
\end{equation}
Si el mismo transductor positivo $B$ satisface
$B(T_H(R))=\tau_R$ y
$B(\Theta_{{\rm clk},a})=\epsilon_a/\ln3$, entonces
\begin{equation}
 \frac{T_H(R)}{\Theta_{{\rm clk},a}}
 =\frac{L_\alpha}{R}\frac{\ln3}{2\pi}.
 \label{eq:ii-ter27-clock-horizon}
\end{equation}
\end{theorem}
\begin{proof}
Restar las esperanzas de \eqref{eq:ii-ter27-modular} da
$\delta s=\Delta\Tr[(\sigma-\rho_A)D]-D(\sigma\Vert\rho_A)$.
Multiplicar por $\tau_R$ prueba \eqref{eq:ii-ter27-free-area}; la
positividad y sus condiciones ya se demostraron en
\ref{thm:ii-area-ley-finita}. Una aplicación lineal positiva entre rectas
preserva cocientes positivos. Por ello el cociente de temperaturas es
$\tau_R/(\epsilon_a/\ln3)$; \eqref{eq:ii-ter27-energia} lo evalúa.
\end{proof}

En $R=L_\alpha$, el factor de conversión es $\ln3/(2\pi)$.
Las dos secciones corresponden a escalas energéticas diferentes y quedan
comparadas por un transporte exacto. Los factores de Barbero convierten
ocupación en área; $\Delta D$ conserva el coste de procedencia; $\tau_R$
lo expresa en energía. La identidad finita conserva además el coste
informacional del estado de ensayo respecto de su referencia.

\subsection{Referencia espectral y multiplicidades conservadas}
\label{subsec:ii-ter27-reference}

El lazo bilateral anterior produce la razón espectral $r=1/10$, y el
agrupamiento trítico conserva $q=r^3=1/1000$ y el resto de cada índice.
El primer retorno estricto al régimen de capacidad, después de la primera
vacancia, fija $C_{24}=23$. El portador mantiene sus tres sectores de
dimensiones $d_0=1$, $d_1=380$, $d_2=649$; sus construcciones de
incidencia, doble graduación y memoria se han demostrado en la sección
precedente. Conservamos la asignación allí declarada de esos sectores a
los grados $23+3n$. Esta asignación es un dato explícito de la realización
conjunta, distinto del cálculo que produce los tres recuentos.

En el espacio $\mathcal D=\bigoplus_{n=0}^2\mathbb C^{d_n}$, sea
\begin{equation}
 \eta=\bigoplus_{n=0}^2r^{23+3n}I_{d_n},\qquad
 b_*:=\Tr\eta
 =10^{-23}\left(1+\frac{380}{1000}+\frac{649}{1000^2}\right).
 \label{eq:ii-ter27-eta}
\end{equation}
La evaluación exacta da $b_*=1.380649\times10^{-23}$ y $0<b_*<1$.
El operador conserva cada etiqueta del portador; la traza reúne sus pesos.
Durante las variaciones termodinámicas siguientes $\eta$ permanece fijo:
es la referencia construida, mientras el estado del sistema de ensayo varía.
Las operaciones, los grados y la normalización quedan así fijados junto
con la referencia que interviene en la variación.

\subsection{Incremento entrópico del registro tensorial}
\label{subsec:ii-ter27-tensor}

Sea $\mathcal H$ un espacio de dimensión finita y $\rho$ una matriz de
densidad sobre él. Escribimos $s(\rho)=-\Tr(\rho\ln\rho)$ y, para
cualquier operador positivo $X$ de traza finita,
$\mathscr H(X)=-\Tr(X\ln X)$, con $0\ln0=0$. La segunda función
admite referencias cuya traza difiere de uno. El registro de ensayo
$\eta\otimes\rho$ conserva por separado la referencia y el sistema.

\begin{theorem}[Entropía marcada del registro producto]
\label{thm:ii-ter27-tensor}
Se cumple exactamente
\begin{equation}
 \mathscr H(\eta\otimes\rho)-\mathscr H(\eta)=b_*s(\rho).
 \label{eq:ii-ter27-tensor}
\end{equation}
\end{theorem}
\begin{proof}
Sean $a_i$ y $p_j$ los autovalores respectivos. Los del producto son
$a_ip_j$, de modo que
\[
 -\sum_{i,j}a_ip_j\ln(a_ip_j)
 =-\sum_i a_i\ln a_i\sum_jp_j
   -\sum_i a_i\sum_jp_j\ln p_j.
\]
Las identidades $\sum_jp_j=1$ y $\sum_i a_i=b_*$ prueban la fórmula.
La convención $0\ln0=0$ incluye estados de rango incompleto por continuidad.
\end{proof}

La diferencia separa la contribución informacional de la referencia en
un balance registrado. La operación tensorial especifica la independencia
del ensayo respecto de esa referencia; un estado correlacionado requiere
su correspondiente información condicional. El coeficiente $b_*$ aparece
por factorización de la traza del registro construido.

\subsection{Realización normalizada y lectura energética condicionada}
\label{subsec:ii-ter27-flag}

La referencia admite una ampliación normalizada en
$\widehat{\mathcal D}=\mathcal D\oplus\mathbb C f$:
\begin{equation}
 \widehat\eta=\eta\oplus(1-b_*)|f\rangle\langle f|,
 \qquad \Omega_\rho=\widehat\eta\otimes\rho.
 \label{eq:ii-ter27-flag}
\end{equation}
Sea $P$ el proyector al sumando $\mathcal D$, extendido por la identidad
de $\mathcal H$, y sea $H=H^*$ el operador energético del sistema,
con origen registrado. La marca $f$ completa la traza de esta realización
auxiliar. El archivo TPK recibido, con sus rutas y memoria, conserva su
propio factor; se transporta junto al registro y no se sustituye por $f$.

\begin{proposition}[Dos funcionales de un estado normalizado]
\label{prop:ii-ter27-readers}
Con los logaritmos restringidos al soporte, las lecturas
\begin{align}
 \mathcal S_P(\rho)&=-\Tr[P\Omega_\rho(I\otimes\ln\rho)]
                    =b_*s(\rho),\label{eq:ii-ter27-S}\\
 \mathcal E_P(\rho)&=
 \frac{\Tr[P\Omega_\rho(I\otimes H)]}{\Tr(P\Omega_\rho)}
                    =\Tr(\rho H)\label{eq:ii-ter27-E}
\end{align}
preservan respectivamente la información relativa a la marca y la energía
por estado condicionado a ella. El denominador es $\Tr(P\Omega_\rho)=b_*$.
\end{proposition}
\begin{proof}
Sobre el sumando marcado, $P\Omega_\rho=\eta\otimes\rho$.
La factorización de la traza da $b_*s(\rho)$ y
$b_*\Tr(\rho H)$ en los dos numeradores. Dividir el segundo por
$b_*>0$ prueba la lectura energética. Tanto $\eta$ como $\rho$ se
recuperan del estado ampliado y del marcador: la primera por su reducción
marcada, y la segunda por la traza sobre $\widehat{\mathcal D}$.
\end{proof}

El bloque complementario contiene también $\rho$. La lectura informacional
total satisface $s(\Omega_\rho)-s(\widehat\eta)=s(\rho)$ y su
energía media total es $\Tr(\rho H)$. Por tanto, elegir el funcional
marcado y la energía condicionada especifica dos observaciones de un estado
que conserva ambos bloques. La elección de observables y la eliminación
física de información son operaciones diferentes.

\subsection{Transductor fijado por la regla conjunta de lectura}
\label{subsec:ii-ter27-transducer}

Trabajamos en las bases positivas transportadas de las rectas de energía,
temperatura y entropía, con $u_S=u_E/u_\Theta$. En las fórmulas de
coordenadas, $H$ expresa energía en $u_E$ y $\mathcal S_P$ expresa
entropía en $u_S$. La regla de realización consta de cuatro operaciones:
recibir la referencia fija \eqref{eq:ii-ter27-eta}; evaluar
\eqref{eq:ii-ter27-S}; evaluar \eqref{eq:ii-ter27-E}; y definir la
sección térmica como conjugada de esos dos funcionales. Las bases se
transportan con los operadores, sin introducir una temperatura objetivo.

\begin{theorem}[Selección térmica dentro de la realización marcada]
\label{thm:ii-ter27-transducer}
Sea $H$ no escalar y sea
$\rho_\beta=e^{-\beta H}/Z(\beta)$ para $\beta>0$. La regla anterior
produce
\begin{equation}
 \frac{d\mathcal S_P}{d\mathcal E_P}=b_*\beta,
 \qquad \Theta_P=\frac1{b_*\beta},\qquad
 B_*(\Theta_Pu_\Theta)=\beta^{-1}u_E,
 \label{eq:ii-ter27-transducer}
\end{equation}
donde $B_*(u_\Theta)=b_*u_E$ es una aplicación lineal positiva única.
\end{theorem}
\begin{proof}
Con $E=\Tr(\rho_\beta H)$, la diferenciación de la suma espectral
finita da $E'=-\operatorname{Var}_{\rho_\beta}(H)<0$ y
$(\ln Z)'=-E$. De $s(\rho_\beta)=\beta E+\ln Z$ se deduce
$s'=\beta E'$. La referencia permanece fija, por lo que
$\mathcal S_P'=b_*s'$ y $\mathcal E_P'=E'$. Su cociente y su
inverso dan las dos primeras identidades. La imagen de una base positiva
determina una única aplicación lineal; sustituir $\Theta_P$ prueba la última.
\end{proof}

El factor depende de la pareja concreta de lecturas. Si se utilizan
simultáneamente energía y entropía marcadas, $(b_*E,b_*s)$, o ambas
condicionadas, $(E,s)$, la derivada entrópica respecto de la energía es
$\beta$. En la pareja \eqref{eq:ii-ter27-S}--\eqref{eq:ii-ter27-E}
es $b_*\beta$. Estas tres identidades se comprueban por la misma derivación;
explican qué normalización determina el coeficiente. Una vez fijada la
regla marcada, multiplicar sólo su canal entrópico por otro factor cambia
la regla y deja de preservar sus hipótesis.

\begin{corollary}[Principio variacional marcado]
\label{cor:ii-ter27-variational}
Para $\Theta>0$, el funcional
\[
 \Phi_\Theta(\rho)=\mathcal E_P(\rho)-\Theta\mathcal S_P(\rho)
 =\Tr(\rho H)-b_*\Theta s(\rho)
\]
posee un único mínimo sobre las matrices de densidad:
\begin{equation}
 \rho_\Theta^*=
 \frac{\exp[-H/(b_*\Theta)]}{\Tr\exp[-H/(b_*\Theta)]}.
 \label{eq:ii-ter27-minimum}
\end{equation}
\end{corollary}
\begin{proof}
Sustituir $\ln\rho_\Theta^*=-H/(b_*\Theta)-\ln Z$ da
\[
 \Phi_\Theta(\rho)-\Phi_\Theta(\rho_\Theta^*)
 =b_*\Theta D(\rho\Vert\rho_\Theta^*)\ge0.
\]
La referencia tiene rango completo. La igualdad en la positividad de la
entropía relativa ocurre exactamente para $\rho=\rho_\Theta^*$,
lo que prueba existencia y unicidad del mínimo.
\end{proof}

La referencia fija y el par de funcionales
\eqref{eq:ii-ter27-S}--\eqref{eq:ii-ter27-E} son las condiciones
constitutivas de esta realización. Bajo ellas, la traza $b_*$ determina
el cociente termométrico y el estado de equilibrio; la identidad de
entropía relativa demuestra la existencia y unicidad del mínimo.

\subsection{Composición con la sección cronológica y sus transportes}
\label{subsec:ii-ter27-map}

Consideremos la realización cronológica de capacidad sobre
$\mathcal D=\bigoplus_{n=0}^2\mathbb C^{d_n}$, con
$(d_0,d_1,d_2)=(1,380,649)$ y $N|_{\mathbb C^{d_n}}=nI$.
Conservamos la asignación de grados $C=23I+3N$ de la sección
\ref{sec:ii-ter27-equivalencia}. El operador de
\eqref{eq:ii-ter27-hcap} es
\begin{equation}
 H_{{\rm clk},a}:=H_{\rm cap}
 =\epsilon_a\frac{\ln10}{\ln3}(23I+3N),\qquad
 \beta_{{\rm clk},a}=\frac{\ln3}{\epsilon_a}.
 \label{eq:ii-ter27-clock-operator}
\end{equation}
Es autoadjunto y positivo en este espacio finito. Sus tres autovalores
distintos, proporcionales a $23,26,29$, prueban que es no escalar.
La multiplicación y el cálculo espectral dan exactamente
\begin{align}
 \beta_{{\rm clk},a}H_{{\rm clk},a}
   &=(23I+3N)\ln10,\nonumber\\
 e^{-\beta_{{\rm clk},a}H_{{\rm clk},a}}
   &=10^{-23}q^N=\eta,\qquad q=10^{-3},\nonumber\\
 \rho_{{\rm clk},a}
   &=\frac{q^N}{Z_N(q)}=\frac{\eta}{b_*},\qquad
 Z_N(q)=1+380q+649q^2.
 \label{eq:ii-ter27-clock-state}
\end{align}
El factor de origen $10^{-23}$ se conserva en la traza $b_*$ y cancela
al normalizar el estado. Para aplicar el teorema
\ref{thm:ii-ter27-transducer}, se fijan la orientación y $\epsilon_a$,
y por tanto $H_{\rm cap}$; la derivada se toma respecto de $\beta$ y
se evalúa después en $\beta_{{\rm clk},a}$. La referencia $\eta$,
independiente de $\epsilon_a$, permanece fija en el primer factor,
mientras $\rho_\beta$ varía en una segunda copia de $\mathcal D$ con
Hamiltoniano fijo \eqref{eq:ii-ter27-clock-operator}. Su evaluación en
$\beta=\beta_{{\rm clk},a}$ proporciona el estado
\eqref{eq:ii-ter27-clock-state} y la energía por nat
$\tau_{{\rm clk},a}=\epsilon_a/\ln3$.

La sección relativa de Williamson mantiene
$H_{W,\mathrm{rel}}=\epsilon_a C$ y
$\tau_W=\epsilon_a/\ln10$. Por tanto,
$H_{{\rm clk},a}/\tau_{{\rm clk},a}
 =H_{W,\mathrm{rel}}/\tau_W=C\ln10$:
coinciden los exponentes y el estado, y las escalas se relacionan por
$\tau_{{\rm clk},a}/\tau_W=\ln10/\ln3$.
La regla marcada, con sus dos funcionales fijados, produce
\begin{equation}
 \Theta_{{\rm clk},P,a}=\frac{\epsilon_a}{b_*\ln3},\qquad
 B_*(\Theta_{{\rm clk},P,a}u_\Theta)
 =\frac{h_a}{108t_0\ln3}u_E.
 \label{eq:ii-ter27-section-test}
\end{equation}
La temperatura se obtiene aquí de la derivada entrópica probada, no de
usar el resultado numérico como condición de ajuste.

Sean $J_E$ el transporte efectivo de energía a estas coordenadas y
$J_\Theta$ el transporte entre las secciones termométricas del mismo
estado de referencia. En particular, $J_\Theta$ envía la sección
cronológica anterior a la sección de \eqref{eq:ii-ter27-section-test},
según sus dos lectores de temperatura. La prueba de unicidad entre rectas
proporciona el cuadrado
\begin{equation}
 B_*J_\Theta=J_EB_{\rm clk}.
 \label{eq:ii-ter27-intertwiner}
\end{equation}
En efecto, las dos aplicaciones coinciden en la sección positiva del
reloj por \eqref{eq:ii-ter27-section-test}, y esa sección genera la recta.
El transporte declara las dos realizaciones comparadas: la igualdad
tipada conserva sus bases y sus secciones; no identifica por simple
notación dos coordenadas térmicas fijadas de manera distinta.

Para el horizonte se obtiene
$\Theta_{H,P}(R)=\tau_R/b_*$. Sustituir en
\eqref{eq:ii-ter27-clock-horizon} mantiene exactamente el cociente entre
ambas temperaturas. Además, con $S_P=b_*s$, la identidad
\eqref{eq:ii-ter27-free-area} toma la forma
\[
 \delta E_R-\Theta_{H,P}(R)\,\delta S_P
 =\tau_RD(\sigma\Vert\rho_A).
\]
Esta composición incorpora el nuevo lector a un balance previo sin
alterar la incidencia, la acción, el área ni el término de entropía relativa.

\subsection{Origen energético, reciprocidad y sección electrónica}
\label{subsec:ii-ter27-electron}

La referencia energética interviene también en las aplicaciones radiales.
Para un carácter positivo $Y$ transportado, se conservan
\[
 H_a=\epsilon_aY,\quad M_a=H_a/c^2,\quad
 R_a=L_\alpha Y,\quad \Lambda_a=L_\alpha Y^{-1}.
\]
Así $H_a\Lambda_a=\hbar_acI$, $R_a\Lambda_a=L_\alpha^2I$ y
$R_a=(G_a/c^2)M_a$. Un desplazamiento $Y\mapsto Y+\delta I$,
con positividad conservada, cambia masa y radio respectivamente en
$\epsilon_a\delta/c^2$ y $L_\alpha\delta$, aunque el desplazamiento
escalar del Hamiltoniano cancele en su estado térmico normalizado.
La procedencia marcada de la energía distingue esas realizaciones.

En el registro de capacidad, el origen recibido es
$Y_{\rm cap}-Y_{\rm rel}=23\ln10/\ln3$.
La diferencia energética es, por tanto,
$\epsilon_a(23\ln10/\ln3)I$; su contribución radial se conserva
en las fórmulas anteriores. Si una distribución se reconstruye desde
$-\ln\rho$, también debe restituirse $\ln Z$ antes de leer su energía
absoluta. El marcador permite conservar ambos datos.

El corrector electrónico ya construido verifica
\begin{align}
 \kappa_e&=80-54\alpha+6\Delta_4,&
 \Delta_4&=\frac{\pi-e\ln\pi}{270}
                 \left(1+\frac\pi{729}\right),\nonumber\\
 \frac{\mathcal E_e^{\beta}}{\mathcal E_e^0}
 &=R_{\rm act}^{\kappa_e},&
 R_{\rm act}&=\frac{\hbar_{\rm pre}}{\hbar_{\rm ret}}.
 \label{eq:ii-ter27-electron-return}
\end{align}
Las energías electrónicas se expresan en la misma unidad y
$\kappa_e>0$ en la sección electrónica construida. El cociente
positivo y su logaritmo fijan la potencia real. Con $L_\alpha$, $c$ y
el transductor comunes, las identidades previas implican
\begin{equation}
 \left(\frac{\mathcal E_e^{\beta}}{\mathcal E_e^0}\right)^{1/\kappa_e}
 =R_{\rm act}
 =\frac{\Theta_{{\rm clk,pre}}}{\Theta_{{\rm clk,ret}}}
 =\frac{G_{\rm ret}}{G_{\rm pre}}.
 \label{eq:ii-ter27-return-ratios}
\end{equation}
Para la rama de retorno, sean $\mathcal E_e=m_ec^2$ y
$y_e=\mathcal E_e/\epsilon_{\rm ret}$. Una sustitución directa produce
\begin{equation}
 \Lambda_e=\frac{\hbar_{\rm ret}c}{\mathcal E_e},\qquad
 \frac{G_{\rm ret}m_e^2}{\hbar_{\rm ret}c}=y_e^2,\qquad
 \beta_{{\rm clk,ret}}\mathcal E_e=(\ln3)y_e.
 \label{eq:ii-ter27-electron-gravity}
\end{equation}
En efecto, $\epsilon_{\rm ret}=\hbar_{\rm ret}c/L_\alpha$ y
$G_{\rm ret}=c^3L_\alpha^2/\hbar_{\rm ret}$ cancelan las escalas
en el segundo cociente, mientras el tercero usa la definición de
$\beta_{{\rm clk,ret}}$. Estas relaciones son composiciones del mismo
retorno, con la conversión dimensional electrónica conservada.

La asignación de grados $C=23I+3N$, la referencia espectral fija $\eta$
y el par de funcionales marcado--condicionado determinan conjuntamente
la realización térmica estudiada. Su prueba tensorial y variacional
establece el coeficiente y el equilibrio; la composición cronológica,
areal y radial conserva el estado, las escalas y el origen energético
que requieren las lecturas respectivas.
